\chapter{Abstract}
\begin{SingleSpace}
This diploma project presents the design, implementation, and experimental evaluation of \textbf{ReconX}, an external penetration testing framework for web applications. The proposed framework follows a black-box methodology and combines passive reconnaissance, active fingerprinting, configuration analysis, and a wide spectrum of vulnerability probes against a remote target whose internal source code and credentials are not available to the tester.

The framework is implemented in \textbf{Python 3} and is organized as a pipeline of 46 specialized modules grouped into 15 execution stages. Reconnaissance modules gather subdomain, DNS, port, and technology information; configuration modules audit \textbf{TLS}, security headers, \textbf{CORS}, authentication cookies, and \textbf{JWT} tokens; active probes test for the most relevant \textbf{OWASP Top 10} classes, including \textit{reflected XSS}, \textit{SQL injection}, \textit{IDOR}, \textit{host header injection}, \textit{cache poisoning}, \textit{prototype pollution}, \textit{open redirect}, \textit{SSRF/SSTI}, \textit{XXE}, \textit{HTTP request smuggling}, and \textit{race conditions}. A central \textit{finding registry} normalizes every output into a unified schema with \textbf{CWE}, \textbf{CVSS}, OWASP, and MITRE ATT\&CK metadata; a correlator and asset-risk module then prioritize the most exploitable assets. A local \textbf{Large Language Model} (Ollama / DeepSeek-R1) generates an explanatory analytical report at the end of each scan.

Experimental evaluation was carried out in two complementary settings: a controlled laboratory run against the deliberately vulnerable \textit{OWASP Juice Shop} target, and an authorized real-world case study against the \texttt{miras.app} academic management system, scanned in the framework's read-only \texttt{bug\_bounty} profile. The framework completed the real-target pipeline in approximately 1\,m~43\,s and produced 721 distinct findings spread across the OWASP Top~10 categories, ranked the 49 discovered assets into four risk tiers (with 5 \textit{critical}-tier assets), and emitted two synthetic cross-probe findings through the correlator. The detected vulnerabilities matched the expected categories of both test targets, only one of 4\,227 issued HTTP requests was outside the configured scope (and that request was blocked by the framework before being sent), and the AI-generated executive summary correctly highlighted the highest-risk asset at the top of the report.

The results confirm that the proposed method makes external web-application penetration testing more reproducible, more explainable, and significantly more time-efficient than a fully manual approach, while keeping the operator in control of intrusive actions through an explicit profile system (\textit{safe}, \textit{bug\_bounty}, \textit{deep}, \textit{intrusive}).

\keywords{external penetration testing; web application security; OWASP Top 10; reconnaissance; vulnerability scanning; finding correlation; LLM-assisted security analysis; ReconX.}
\end{SingleSpace}
